% TOP_PROBLEM: {"problemId": "problem.top500-omission-w041044-n-el-order-for-the-spin-1-2-heisenberg-antiferromagnet-on-the-square-lattice", "canonicalTitle": "Néel Order for the Spin-1/2 Square-Lattice Heisenberg Antiferromagnet (Averaged Ground-State Order)", "releaseRank": 325, "primaryDomain": "mathematical_physics", "primaryDomainLabel": "Mathematical physics", "displayStatus": "open", "releaseStatus": "open"}
% !TeX program = lualatex
% ATTEMPT_STATUS: unresolved
% Model: GPT-6 Astra Ultra; continuation dated 27 September 2026.
\documentclass[11pt]{article}
\usepackage[margin=25mm]{geometry}
\usepackage{fontspec}
\setmainfont{Latin Modern Roman}
\usepackage{amsmath,amssymb,mathtools}
\usepackage{unicode-math}
\setmathfont{Latin Modern Math}
\usepackage{xurl,hyperref}
\hypersetup{colorlinks=true,urlcolor=blue}
\setcounter{secnumdepth}{0}
\setlength{\parindent}{0pt}
\setlength{\parskip}{5pt}
\setlength{\emergencystretch}{3em}
\begin{document}
\section*{\texorpdfstring{Néel Order for the Spin-1/2 Square-Lattice Heisenberg Antiferromagnet (Averaged Ground-State Order)}{Néel Order for the Spin-1/2 Square-Lattice Heisenberg Antiferromagnet (Averaged Ground-State Order)}}\label{325}
\subsection{Definitions and mathematical statement}
For even $L$, put $\Lambda_L=({\mathbb{Z}}/L{\mathbb{Z}})^2$ and associate ${\mathbb{C}}^2$ to each site. Let $S_x^a=\frac12\sigma_x^a$ be the spin-$\frac12$ operators and
\[
 H_L=\sum_{\langle x,y\rangle}\sum_{a=1}^3S_x^aS_y^a,
 \qquad M_L=\sum_{x\in\Lambda_L}(-1)^{x_1+x_2}S_x^3.
\]
The sum is over nearest-neighbor bonds with periodic boundary conditions and positive coupling normalized to one. For the normalized ground state $\Phi_L$, the selected N\'eel-order assertion is
\[
 \liminf_{\substack{L\to\infty\\L\text{ even}}}
 \langle\Phi_L,(M_L/L^2)^2\Phi_L\rangle>0.
\]
This is a nonzero thermodynamic lower bound on the squared averaged staggered magnetization. It is not a claim that the finite-volume symmetric ground state has nonzero first moment of $M_L$.
\par\smallskip\textit{Formulation and concept references:} \href{https://arxiv.org/pdf/1807.05847}{[S1]}.

\subsection{Short English statement}
Does the spin-one-half square-lattice antiferromagnet retain nonzero squared staggered magnetic order in its ground state as the lattice size tends to infinity?
\subsection{Research attempt}
\textbf{Scope, symmetry, and the source's logical direction.}
Write $N=L^2$ and take even $L\geq4$, so there are $2N$ ordinary
nearest-neighbor bonds without the small-torus convention ambiguity.
The source [S1] explicitly singles out dimension two and spin one half as the
exception to the general rigorous long-range-order results it discusses.
Its tower-of-states conclusions assume long-range order; they cannot be used
in reverse to establish that missing assumption.  A current search also found
numerical work on frustrated $J_1$--$J_2$ and other lattices, whose numerical
phase identifications do not establish this thermodynamic lower bound for the
specified Hamiltonian.

The finite even bipartite model has a unique singlet ground state, by the
Marshall--Lieb--Mattis result recalled in [S1].  We use that established
finite-volume theorem as input.  Spin-rotation invariance then forces
$\langle M_L\rangle=0$, but says nothing by itself about the size of
$\langle M_L^2\rangle$.  In fact $0\leq\langle M_L^2\rangle\leq N^2/4$
because $\|M_L\|\leq N/2$.  The question is whether the ground state retains
a fixed positive fraction of this quadratic scale, rather than only order
$N$ fluctuations.

\textbf{Bond algebra and two competing trial states.}
For two spin-one-half sites,
\[
 S_x\cdot S_y=\frac12 P_{xy}-\frac14 I,
\]
where $P_{xy}$ swaps the two tensor factors.  Its eigenvalues are $-3/4$
on the singlet and $1/4$ on the triplet.  Hence the crude operator lower
bound on the lattice is $H_L\geq-3N/2$.  The alternating product state with
up spins on one sublattice and down spins on the other has each bond
expectation $-1/4$, and therefore
\[
 \langle H_L\rangle_{\rm N\acute eel}=-N/2,
 \qquad \langle M_L^2\rangle_{\rm N\acute eel}=N^2/4.
\]
This supplies a variational upper bound for the ground energy; it does not
compare the trial state's magnetization with the ground state's magnetization.

A contrasting state gives a concrete obstruction to that inference.  Tile
the torus by disjoint $2\times2$ plaquettes.  On one plaquette label its two
diagonal sublattices $A$ and $B$, and set $S_A=S_1+S_3$,
$S_B=S_2+S_4$.  The four cycle bonds form the complete bipartite graph
$K_{2,2}$, so
\[
 H_\square=S_A\cdot S_B
 =\frac12\bigl((S_A+S_B)^2-S_A^2-S_B^2\bigr).
\]
Take the state with $S_A=S_B=1$ coupled to total spin zero.  It has energy
$-2$, and its staggered component is $M_\square=S_A^3-S_B^3$.
On this singlet $S_B^3=-S_A^3$; rotational invariance gives
$\langle(S_A^3)^2\rangle=S_A(S_A+1)/3=2/3$.  Thus
\[
 \langle M_\square^2\rangle=8/3,
 \qquad\langle S_x^a\rangle=0.
\]
Take the tensor product of these plaquette singlets.  Bonds joining different
plaquettes have zero expectation, since their single-spin means factorize
and vanish.  With $N/4$ plaquettes this gives
\[
 \langle H_L\rangle_{\rm plaq}=-N/2,
 \qquad\langle M_L^2\rangle_{\rm plaq}=\frac{2N}{3}.
\]
The squared average tends to zero.  Thus the same elementary energy density
as the ordered product state is compatible with no macroscopic staggered
variance.  Neither state is asserted to be the full lattice ground state.
The example rules out an argument using only that variational energy value;
it does not show that stronger energy/correlation estimates cannot work.

\textbf{An exact Fourier sum rule and the desired concentration.}
For momenta $k\in(2\pi/L)\{0,\ldots,L-1\}^2$, define
\[
 S_k^a=N^{-1/2}\sum_xe^{-ik\cdot x}S_x^a,
 \qquad s_L(k)=\langle S_k^3 S_{-k}^3\rangle\geq0.
\]
Positivity follows from $S_{-k}^3=(S_k^3)^*$; the two operators commute.
Fourier orthogonality gives, in every state,
\[
 \sum_k s_L(k)=\sum_x\langle(S_x^3)^2\rangle=N/4.
\]
At $Q=(\pi,\pi)$, one has $M_L=\sqrt N S_Q^3$, hence
\[
 \frac{\langle M_L^2\rangle}{N^2}=
 \frac{s_L(Q)}{N}=\frac14-\frac1N\sum_{k\ne Q}s_L(k).
\]
The exact question is therefore whether a positive fraction of the sum-rule
mass remains at the antiferromagnetic momentum.  Positivity of correlations
at a finite collection of distances does not automatically imply this atom
of order $N$.

The energy also has an exact momentum representation:
\[
 H_L=\sum_k(\cos k_1+\cos k_2)
                    \sum_{a=1}^3 S_k^aS_{-k}^a.
\]
It follows by Fourier inversion of the two positively oriented bond sums;
the imaginary parts cancel between $k$ and $-k$.  In the singlet ground state,
rotational invariance makes the three components equal, so its expectation is
$3\sum_k(\cos k_1+\cos k_2)s_L(k)$.
This constrains a weighted moment of the structure factor.  Such a moment
and the total-mass sum rule alone do not force macroscopic mass at the single
point $Q$: mass can accumulate in a shrinking neighborhood without an atom
at that exact momentum for finite $L$.  An estimate controlling all the
non-$Q$ modes is needed.

\textbf{A conditional infrared criterion with the constants exposed.}
Let $\varepsilon(k)=2+\cos k_1+\cos k_2$, which vanishes precisely at $Q$.
If one could prove a uniform bound
\[
 s_L(k)\leq\frac{C}{\sqrt{\varepsilon(k)}}\qquad(k\ne Q)
\]
with an explicit constant satisfying
\[
 C I<\frac14,\qquad
 I=\frac1{(2\pi)^2}\int_{[-\pi,\pi]^2}
                \frac{dk}{\sqrt{2+\cos k_1+\cos k_2}},
\]
then the sum rule would prove the requested positive liminf.
The integral is finite: near $Q$, the denominator is comparable to
$|k-Q|$, whose reciprocal is locally integrable in two dimensions.
The corresponding punctured discrete averages converge to $I$.  For a
direct justification, away from a small disk ordinary Riemann sums converge;
inside the disk, grouping lattice points into annuli of grid-width size gives
a bound proportional to the disk radius, uniformly in $L$.  Then let the
radius go to zero.

This criterion deliberately separates integrability from a sufficient
constant.  A finite infrared integral does not establish $CI<1/4$.
At the spin-one-half endpoint, inserting an unproved improved coefficient
would assume the difficult part.  Nor can a positive-temperature estimate
with a different singularity be substituted: a bound proportional to
$1/\varepsilon(k)$ has a logarithmically divergent two-dimensional integral.
The quantum zero-temperature problem requires its own estimates and order
of limits.

\textbf{A rigorous implication from order to a low-energy excitation.}
Let $\Phi$ be the singlet ground state, with energy $E_0$, and assume
$\langle M_L^2\rangle>0$.  Then $M_L\Phi$ is orthogonal to $\Phi$ and its
Rayleigh quotient above the ground energy is
\[
 \frac{\langle\Phi,M_L(H_L-E_0)M_L\Phi\rangle}
      {\langle M_L^2\rangle}
 =\frac{\langle[M_L,[H_L,M_L]]\rangle}
        {2\langle M_L^2\rangle}.
\]
For a nearest-neighbor bond write
$H_{xy}^{\perp}=S_x^1S_y^1+S_x^2S_y^2$.
Since its sites have opposite staggering signs, the spin commutation relations
give
\[
 [M_L,[S_x\cdot S_y,M_L]]=-4H_{xy}^{\perp},
 \qquad \|H_{xy}^{\perp}\|=1/2.
\]
Summing $2N$ bonds bounds the numerator in the Rayleigh quotient by $2N$.
The first excitation gap thus obeys
\[
 E_1-E_0\leq\frac{2N}{\langle M_L^2\rangle}.
\]
If the desired order bound $\langle M_L^2\rangle\geq cN^2$ were known,
this would give a gap at most $2/(cN)$.  This is an elementary first step
toward a tower-of-states consequence.  Reversing the implication is invalid:
the existence of small-energy excitations need not make this particular
matrix element or order parameter large.

\textbf{Verification and outcome.}
The saved diagnostic forms the four-spin plaquette matrices, checks its
ground energy and staggered variance numerically against $-2$ and $8/3$,
and verifies the bond double-commutator identity by matrix arithmetic.
It also checks finite Fourier normalization identities.  No finite-size fit
is used as a proof of a positive thermodynamic liminf.  The new derivations
give explicit trial-state counterchecks, a sufficient infrared criterion,
and the one-way excitation bound.  The missing step is a uniform ground-state
estimate that leaves strictly positive sum-rule mass at $Q$ for spin one half
on this two-dimensional lattice.
\subsection{Sources}
\begin{itemize}
\item[S1]Long-range order, tower of states, and symmetry breaking in lattice quantum systems, arXiv:1807.05847v6.
\url{https://arxiv.org/pdf/1807.05847}.
\textit{Formulation source.}
\end{itemize}
\end{document}
